% ---- Résumé français ----
\section*{Résumé}

Ce rapport présente \textbf{TAC-TIC IntelliPM}, une plateforme web complète de gestion
de projets développée dans le cadre d'un projet de fin d'études au sein de l'entreprise
TAC-TIC (Groupe New Ways). L'application couvre l'intégralité du cycle de vie d'un
projet logiciel : planification, suivi agile par sprints, gestion hiérarchique des tâches,
visualisation Gantt avec calcul du chemin critique, et collaboration en temps réel via
WebSocket.

La valeur différenciante de la plateforme repose sur son moteur d'intelligence
artificielle hybride : un modèle Machine Learning à 22 variables (ensemble
\textit{Gradient Boosting} + \textit{Random Forest}) pour la recommandation
d'assignation de tâches, couplé à un LLM (Llama 3.3 70B via OpenRouter) pour
l'estimation de durée et la détection de risques.

L'architecture repose sur trois couches : un backend \textbf{Laravel~12} (PHP~8.2)
exposant~251 routes REST, sécurisé par Laravel Sanctum et Spatie Permission (RBAC
à deux niveaux) ; une SPA \textbf{React~19} (Vite~7 + TailwindCSS~4) avec
internationalisation complète en français ; et un microservice \textbf{FastAPI}
(Python~3.11) dédié aux modèles ML.

La démarche de développement est agile Scrum, structurée en 5~sprints couvrant :
gestion des utilisateurs, projets et équipes, tâches et Kanban, planification visuelle
et notifications, et intelligence artificielle.

\vspace{0.5em}
\noindent\textbf{Mots-clés :} Gestion de projets, Scrum, Machine Learning, Laravel,
React, Kanban, Gantt, FastAPI, Recommandation, WebSocket.

\vspace{1.5em}
\rule{\linewidth}{0.3pt}

% ---- Abstract anglais ----
\section*{Abstract}

This report presents \textbf{TAC-TIC IntelliPM}, a comprehensive web-based project
management platform developed as a final-year project at TAC-TIC (Groupe New Ways).
The application covers the full software project lifecycle: planning, agile sprint
tracking, hierarchical task management, Gantt visualization with critical path
computation, and real-time collaboration via WebSocket.

The platform's differentiating feature is its hybrid AI engine: a 22-feature Machine
Learning model (Gradient Boosting + Random Forest ensemble) for task assignment
recommendation, combined with an LLM (Llama 3.3 70B via OpenRouter) for duration
estimation and risk detection.

The architecture spans three layers: a \textbf{Laravel~12} (PHP~8.2) backend exposing
251 REST routes, secured by Laravel Sanctum and Spatie Permission (two-level RBAC);
a \textbf{React~19} SPA (Vite~7 + TailwindCSS~4) fully localized in French; and a
dedicated \textbf{FastAPI} (Python~3.11) microservice for ML models.

Development follows Scrum across 5 sprints: user management, projects and teams,
tasks and Kanban, visual planning and notifications, and artificial intelligence.

\vspace{0.5em}
\noindent\textbf{Keywords:} Project Management, Scrum, Machine Learning, Laravel,
React, Kanban, Gantt, FastAPI, Recommendation, WebSocket.
